% https://www.overleaf.com/3852372271rfxvjvrbbzzq#551c1d


\documentclass{article}

% Language setting
% Replace `english' with e.g. `spanish' to change the document language
\usepackage[english]{babel}

% Set page size and margins
% Replace `letterpaper' with `a4paper' for UK/EU standard size
\usepackage[letterpaper,top=2cm,bottom=2cm,left=3cm,right=3cm,marginparwidth=1.75cm]{geometry}

% Useful packages
\usepackage{amsmath}
\usepackage{graphicx}
\usepackage[colorlinks=true, allcolors=blue]{hyperref}
\title{AI and Law: exoskeleton case study}
\date{A.A. 2025/2026}
\author{Luca Barbati (12345678),\\Roberto Lazzarini (4937188),\\  Tommaso Viotti (5048661)}
\begin{document}
\maketitle
\newpage

\section{What are the pros and cons of using this exoskeleton?}
Industrial exoskeletons have the potential to radically change many sectors, for example warehouse management. 
If not regulated, the advantages they bring to the table may be undermined by the risks they bring to the rights of the employees and the company itself.
The main and most obvious advantage is an increased efficiency of the worker wearing the exoskeleton, which is accompanied by: 
\begin{itemize}
    \item - a reduction of injury risks (metti qua un esempio), which could decrease the need for employee turnover
    \item - standardized performances: employee output is derived from skill rather than phisical prowess
    \item 
\end{itemize}
The company image can be improved if these systems are employed in a fair and correct way, in fact the wrong approach to them can create friction between the employer and the employee and several other risk:
\begin{itemize}
    \item - exoskeleton requires skill, so training and a form of certification for the usage is needed
    \item - these systems are expensive and need maintenance done by a qualified employee
    \item - can introduce discrimination, for example for health conditions, as shown in the scenario
    \item - even if it increase efficency it can limit movements and be uncomfortable to wear for a long time, producing a strong stress for the employee
    \item - experience workers may be skeptical towards a different work paradigm (ex. "I have been working with my bare hands for twenty years, what do i need this piece of junk for?")
    \item - data recorded by the system may lead to privacy violation
\end{itemize}

Consequently, there is a trade-off in the usage of these specific systems and well-thought regulations must be enforced.


\section{It is possible to impose the use of these systems?}

Regarding the employer’s ability to require exoskeleton usage, it is necessary to differentiate between two scenarios: employees who were not originally contracted to operate an exoskeleton (scenario 1), and employees who were hired specifically for roles that involve its use (scenario 2).

The most crucial aspects of the first scenario are consent and negotiation: only take on a contractually unforeseen task if they agree to it.
In that case, they must be adequately trained and their pay grade may be revised through negotiation, potentially involving a trade union, to level out contractual power.  

In the second scenario the employer is specifically hiring someone to work with an exoskeleton. In such a scenario, it is acceptable to demand the ability to operate with an exoskeleton, just like it is possible to demand someone with a driving license. 

In both scenarios the medical condition of the employee must be evaluated before working with an exoskeleton. 
Moreover, in case a debilitating medical condition arises after hiring or the task change process, the employer must move the employee to a different task (at the same pay grade). Only in extreme cases an employee can be fired due to medical condition. The trade union must be involved in both processes. 

\section{How could they be introduced and regulated?}
In terms of regulations and introduction of the exoskeleton several aspects must be take into account.
Starting from the producer, it must provide extensive testing to ensure high security standard for user and its surroundings. Furthermore, it must not collect personalized user data but it should be done by the company and specifically the safety department, where the safety manager is the only figure allowed to access user safety data.

- producer must provide extensive testing to ensure high security standard for user and its surroundings
- Producer must not collect personalized user data 

Classification
\begin{itemize}
    \item - Each exoskeleton must be appropriately certified and properly classified to adhere to security standards decided by the lawmakers.
    \item - Exoskeletons must be classified by weight class, skill required and type of task.
    \item -- Every exoskeleton belonging to a certain task needs to adhere to the same interface standards to achieve interface uniformity (ex. if you know how to drive a fiat panda you know how to drive an audi)
    \item - To use each exoskeleton class a different licence/certification must be required (as we already do with trucks or forklift and other industrial equipment).
\end{itemize}


Company 
\begin{itemize}
    \item - only certified employees must be allowed to pilot an exoskeleton
    \item - the training must be provided by the company and it must not be enforced
    \item - the employee must consent for acquisition of data
    \item - employee data collection can only be performed for employee safety only during work hours
    \item - employee data cannot be used to evaluate performance and to decide if lay someone down
    \item - the security manager is the only figure allowed to access user security data
\end{itemize}
Company doctor
- exoskeleton are devices particularly integrated with the body of each user. 
- Hence, the mutua's doctor must release certification to use the exoskeleton for each worker and it must be possible edit exoskeleton security parameteres according to each worker
- At this point, in the provider scenario, we can proceed only after the doctor’s permission. After the visit, it is possible to modify Laura’s parameters according to the new range and vital signs defined by the doctor.

Employee
- The employee must follow all the securty rules and to use the exoscheleton (and recive the adequate pay) must agree to the condition describe before.


After all of this the scenario is not possible anymore. Even in the case of a health condition so severe that the doctor doesn t allows the employee to continue working, the employee must be moved to another position and only as a last option can be fired. In these extreme cases the trade union must negotiate the best solution for both the company and the employee.



%- OFFICAL DISCOURSE STARTS HERE
Introducing and regulating the use of exoskeletons requires a structured approach, starting from the responsibilities of the producer. 
Manufacturers must conduct extensive testing to ensure high safety standards for both the users and their surroundings, while refraining from collecting personalized user data. 

Lawmakers must define an appropriate exoskeleton classification to guarantee safety, interoperability, and regulatory compliance: each device should be certified and categorized according to factors such as weight, required skill level, and intended task. Within each category, standardized interfaces should be adopted to ensure consistency and facilitate ease of use, allowing operators to transfer their skills across similar systems. Moreover, each class of exoskeleton should require a specific license or certification, in line with existing regulations for industrial machinery, thereby ensuring that only adequately trained individuals can operate them.

As stated before, at the company level it must be enforced that only certified employees should be allowed to use such devices. Untrained employees tasked with exoskeleton-related duties must be provided proper training.
Exoskeleton are closely integrated with the user's body, so the company doctor plays a crucial role: exoskeleton usage cannot start without medical cleareance and the company doctor must fine-tune and update the exoskeleton’s safety and comfort parameters as needed, to fit the employee.
Employees must give explicit consent for any data collection, which should be strictly limited to safety purposes during working hours and must never be used for performance evaluation or employment decisions. 
Moreover, data collection should be managed internally and only by an appositely appointed safety department and access to sensitive information should be strictly limited to a designated safety manager. 

In certain circumstances, workers could no longer able to perform their tasks. In cases where a worker’s health condition prevents further exoskeleton use, the company is required to reassign the employee to a different role whenever possible, with dismissal considered only as a last resort. In such critical situations, the involvement of trade unions is essential to negotiate a fair and balanced solution that safeguards both the employee’s rights and the company’s needs. 
Overall, this structured system ensures that the adoption of exoskeletons remains safe, ethical, and socially sustainable, clearly defining responsibilities while prioritizing worker protection and regulatory compliance.
% --------------






\section{Who are the stakeholders involved?}
- Lawmaker
- Producer
- Security manager (analogous to IT manager)
- Company
- Company doctor
- Trade union
- Employee


\end{document}
